\documentclass[10pt,a4paper]{amsart}
\usepackage[margin=19mm]{geometry}
\usepackage{amsmath,amssymb,mathtools}
\usepackage[scr=rsfs,cal=euler]{mathalfa}
\usepackage{xfrac}
\usepackage[unicode,hidelinks,bookmarks=false]{hyperref}
\newcommand{\ie}{i.e.\ }
\newcommand{\cf}{cf.\ }
\newcommand{\resp}{respectively\ }
\newcommand{\Subsection}{\subsection}
\providecommand{\nobreakdash}{\nobreak\hbox{-}}
\setlength{\emergencystretch}{2em}
\multlinegap0pt
\usepackage{amssymb}
\usepackage[normalem]{ulem}
\newtheorem{thm}{Theorem}
\newcommand{\Discr}{\mathbb{D}_m}
\newcommand{\R}{\mathbb{R}}
\newcommand{\Z}{\mathbb{Z}}
\newcommand{\eb}{\mathbf{e}} 
\newcommand{\var}{s} 
\title{Harish-Chandra's admissibility problem \\ for Banach space representations of $\mathrm{SL}(2,\R)$}
\author{Francesca Astengo, Michael G. Cowling, Bianca Di Blasio}
\date{Source: arXiv:2406.12626v3 (CC BY 4.0)}
\begin{document}
\maketitle

\noindent\emph{Source and licence.} Harish-Chandra's admissibility problem \\ for Banach space representations of $\mathrm{SL}(2,\R)$. Version arXiv:2406.12626v3 (CC BY 4.0), distributed by arXiv under the Creative Commons Attribution 4.0 International licence, http://creativecommons.org/licenses/by/4.0/, as recorded in the arXiv licence metadata for this exact version and shown on the abstract page https://arxiv.org/abs/2406.12626v3. Full licence notice: \texttt{licenses/arxiv-2406.12626-CC-BY-4.0.txt}.\par\smallskip

\noindent\textit{Editorial scope.} Counted unit: the article's thm thm:f-star-g (source0.tex lines 731--742). The article's complete original proof of it is delivered with it (source0.tex lines 744--846, one \texttt{proof} environment of 103 lines). No mathematics is written, repaired or re-derived: the statement and the proof are the article's own lines, in the article's own notation, and no retained line is substituted or re-flowed.  Retained uncounted as the source-local context the proof needs: lines 436--444; lines 622--624; lines 694--700; lines 696--698.  Nothing was omitted from any retained span, so the package carries no omission record.  No retained line contains a citation, so the wrapper carries no bibliography and no bibliography entry.  No retained line contains a figure environment, an image-inclusion command or drawing code, so no image is delivered and the package declares no asset dependency.  Editorial formatting only, in addition: an AMS \texttt{amsart} preamble that restates the article's own declarations and macros used by the retained lines, namely the environments \texttt{thm} on separate counters, together with the standard packages those lines need.  The retained environments print in this document as Theorem 1 in order of appearance, whereas the article prints the same items with the numbering of its own full document.  The complete native source of the article as distributed in the arXiv e-print for this exact version is bundled unchanged as \texttt{sources/160965/source0.tex}.
\begin{equation}\label{e:plancherel}
\begin{aligned}
\|f\|_{L^2(G)}^2&=\frac{1}{8\pi} \int_\R \|\pi_{i\lambda,0}(f)\|_{HS}^2 \lambda\tanh(\pi\lambda)\, d\lambda
+
 \frac{1}{8\pi} \int_\R \|\pi_{i\lambda,1}(f)\|_{HS}^2 \lambda\coth(\pi\lambda)\, d\lambda
 \\
 &+\frac{1}{8\pi} \sum_{s\in \frac12 \Z \setminus\{0\}} |s| \left(\|\pi_s^+(f)\|_{HS}^2 +\|\pi_s^-(f)\|_{HS}^2 \right) ,
\end{aligned}
\end{equation}

\begin{equation}\label{e:sfer-repr}
\phi_{m,\var}(x)=\langle \pi(x)\eb^{\pi}_m,\eb^{\pi}_m\rangle ,
\end{equation}

\begin{thm}\label{thm:PW}
The (generalised) $m$-spherical transform is bijective  from  $C^\infty_c(G)_m$ onto the space of even entire functions $\psi$ for which there exists a positive constant $C_\psi$ such that
\begin{equation}\label{e:uet}
\sup_{\var\in \mathbb C} (1+|\var|)^{N}\, |\psi(\var)| e^{-C_\psi |\mathrm{Re} \var|}<\infty.
\end{equation}
 for every positive integer $N$.
\end{thm}

\begin{equation}
\sup_{\var\in \mathbb C} (1+|\var|)^{N}\, |\psi(\var)| e^{-C_\psi |\mathrm{Re} \var|}<\infty.
\end{equation}

\begin{thm}\label{thm:f-star-g}
Fix $m,n \in \Z$ such that $m-n$ is even.
Suppose that $g \in C^\infty_c(G)$ and 
\begin{equation}\label{eq:invariance-mn}
g(k_\theta x k_\phi) = e^{-i n\theta-im\phi}  g(x) 
\qquad\forall x \in G  \quad\forall \theta, \phi \in \R.
\end{equation}
Then for all $f \in C^\infty_c(G)_n$, there exists $f' \in C^\infty_c(G)_m$ such that 
\[
f * g = g * f'  .
\]
\end{thm}

\begin{proof}
%%%%%%%%%%%%% INIZIO NOTA
%\footnote{Since  $f \in C^\infty_c(G)_n$, its $n$-spherical transform $\psi$ is an even entire function satisfying~\eqref{e:uet}. Then by Theorem~\ref{thm:PW}, the function $f'$, whose $m$-spherical transform is $\psi$,
%is in $C^\infty_c(G)_m$. 
%By formula~\eqref{e:sfer-repr} this implies that, if 
%$\pi $ is in  the unitary principal series $\left\{\pi_{\var,\epsilon}\right\}_{s\in i\R}$ 
%or the discrete series $\left\{\pi^\pm_{\var}\right\}_{s\in {\Discr}\cap{\mathbb{D}_n}}$,  
%\[
%\langle \pi(f')\eb^\pi_m, \eb^\pi_m\rangle=
%\langle \pi(f)\eb^\pi_n, \eb^\pi_n\rangle.
%\]
%
%%In unitary representation terms, 
%%for $s\in \Discr\cap \mathbb{D}_n$ or $s\in i\R$
%%\footnote{Qui per dopo ho un piccolo problema su cui pensare $\Discr$ e $\mathbb{D}_n$, probabilmente ci vogliono tutti i semiinteri o gli interi... Ecco: bastano quelli in $\Discr\cap\mathbb{D}_n$; per questi i coefficienti diagonali della discreta e della serie principale sono uguali.
%%}
%%\[
%%\langle \pi_{\var,\epsilon}(f)\eb_n,\eb_n\rangle=\psi(\var)=\langle \pi_{\var,\epsilon}(f')\eb_m,\eb_m\rangle .
%%\]
%
%By~\eqref{eq:invariance-mn} if $\pi$ is a unitary representation, then
%\[
%\langle \pi(g)\eb^\pi_h,\eb^\pi_j\rangle=\int_G g(x)\, \langle \pi(x)\eb^\pi_h,\eb^\pi_j\rangle\, dx=0
%\]
%unless  $j=n$ and $h=m$, and the same holds for $f*g$ and $g*f'$, which share the same kind of invariance.
%Therefore we expect a contribution  from the representations of the unitary principal series or the discrete series only with $s\in \Discr\cap{\mathbb{D}_n}$.
%Since
%\begin{align*}
%\langle \pi(f*g)\eb^\pi_m,\eb^\pi_n\rangle=\langle \pi(g)\eb^\pi_m,\pi(f)^*\eb^\pi_n\rangle
%=\langle \pi(g)\eb^\pi_m,\eb^\pi_n\rangle\, \langle \pi(f)\eb^\pi_n, \eb^\pi_n\rangle
%\end{align*}
%and
%\begin{align*}
%\langle \pi(g*f')\eb^\pi_m,\eb^\pi_n\rangle=\langle \pi(f')\eb^\pi_m,\pi(g)^*\eb^\pi_n\rangle
%=\langle \pi(g)\eb^\pi_m,\eb^\pi_n\rangle\, \langle \pi(f')\eb^\pi_m, \eb^\pi_m\rangle,
%\end{align*}
%the thesis follows from the Plancherel formula~\eqref{e:plancherel}.
%}
%%%%%%%%%%%%%%%%%%%%%%%%% FINE NOTA
By~\eqref{eq:invariance-mn}, if $\pi$ is a unitary representation, then
\[
\langle \pi(g)\eb^\pi_h,\eb^\pi_j\rangle=\int_G g(x)\, \langle \pi(x)\eb^\pi_h,\eb^\pi_j\rangle\, dx=0
\]
unless  $j=n$ and $h=m$. Moreover, 
%\[
%\langle \pi(g)\eb^\pi_m,\eb^\pi_n\rangle=0 \quad
%\text{if}\quad
%\begin{cases}
%\pi=\pi_s^+ \text{ and } 2|s|>\min\{m,n\}-1
%\\[3pt]
%\pi=\pi_s^- \text{ and } 2|s|>\min\{-m,-n\}-1.
%\end{cases} 
%\]
%
\[
\langle \pi(g)\eb^\pi_m,\eb^\pi_n\rangle=0 \quad
\text{if}\quad
\begin{cases}
\pi=\pi_\var^+ \text{ and } 2|\var|\not\in \min\{m,n\}-1-2\mathbb N
\\[3pt]
\pi=\pi_\var^- \text{ and } 2|\var|\not\in\min\{-m,-n\}-1-2\mathbb N,
\end{cases} 
\]
so that, in the Plancherel formula of a function 
satisfying~\eqref{eq:invariance-mn}, only a finite number of half-integers $\var$ appear.

As  $f \in C^\infty_c(G)_n$, its $n$-spherical transform $\psi$ is an even entire function satisfying~\eqref{e:uet}. 
By Theorem~\ref{thm:PW}, the function $f'$ whose $m$-spherical transform is $\psi$
is in $C^\infty_c(G)_m$. 

We now check that $f*g=g*f'$. By the Plancherel formula~\eqref{e:plancherel} and 
since  $f*g$ and $g*f'$ both satisfy~\eqref{eq:invariance-mn}, it is enough to check that
$\langle \pi(f*g)\eb^\pi_m,\eb^\pi_n\rangle=\langle \pi(g*f')\eb^\pi_m,\eb^\pi_n\rangle$
in the following cases:
\begin{itemize}
\item $\pi=\pi_{i\lambda,\epsilon}$, where $\lambda\in \R $
and $m-\epsilon \in 2\Z$ (hence $n-\epsilon$ is even);
\item 
$\pi=\pi^+_s$, where $2|s|\in \min\{m,n\}-1-2\mathbb N$, that is $m,n>0$ and $s\in \Discr\cap{\mathbb{D}_n}$;
\item $\pi=\pi^-_s$, where $2|s| \in\min\{-m,-n\}-1-2\mathbb N$, that is $m,n<0$ and $s\in \Discr\cap{\mathbb{D}_n}$.
\end{itemize}
%%%%%%%%%%%%%%%inizio nota 

By formula~\eqref{e:sfer-repr}, 
in all these cases,  \[
\langle \pi(f')\eb^\pi_m, \eb^\pi_m\rangle=
\langle \pi(f)\eb^\pi_n, \eb^\pi_n\rangle,
\]
so that 
\begin{align*}
\langle \pi(f*g)\eb^\pi_m,\eb^\pi_n\rangle
&=\langle \pi(g)\eb^\pi_m,\pi(f)^*\eb^\pi_n\rangle
\\
&=\langle \pi(g)\eb^\pi_m,\eb^\pi_n\rangle\, \langle \pi(f)\eb^\pi_n, \eb^\pi_n\rangle
\\
&=\langle \pi(g)\eb^\pi_m,\eb^\pi_n\rangle\, \langle \pi(f')\eb^\pi_m, \eb^\pi_m\rangle 
\\
&=\langle \pi(f')\eb^\pi_m,\pi(g)^*\eb^\pi_n\rangle
\\
&=\langle \pi(g*f')\eb^\pi_m,\eb^\pi_n\rangle ,
\end{align*}
as required.
\end{proof}

\end{document}
